% =====================================================================
%  Rapport de Projet de Fin d'Études — SUP'COM
%  Sujet  : labConnect — Plateforme digitale de gestion de laboratoire
%           d'analyses médicales (Mauritanie)
%
%  Mise en forme conforme au Guide ID.145 v2 (mai 2026) :
%   - Page de garde standard SUP'COM
%   - Polices limitées, harmonie graphique (Latin Modern)
%   - Interligne 1.3, marges adaptées à la reliure
%   - Texte justifié
%   - Figures et tableaux numérotés et titrés
%   - Références bibliographiques au format [N]
%
%  Compilation : pdflatex main.tex (×2, pour la table des matières)
%                + bibtex si on passe à BibLaTeX (pas requis ici)
% =====================================================================

\documentclass[12pt, a4paper, oneside]{report}

% ---------- Langue, encodage ------------------------------------------
\usepackage[utf8]{inputenc}
\usepackage[T1]{fontenc}
\usepackage[french]{babel}
\frenchsetup{StandardLayout=true, ItemLabeli=\textendash, ItemLabelii=\textendash}

% ---------- Polices ---------------------------------------------------
% Latin Modern : neutre, très lisible, large support des accents français.
% On reste sur une police principale + une mono pour les blocs de code,
% conformément au § 1 du guide (« choix limité, harmonie graphique »).
\usepackage{lmodern}
\usepackage{microtype}      % justifications plus propres

% ---------- Mise en page ----------------------------------------------
\usepackage[
  a4paper,
  inner=2.8cm,   % marge intérieure (côté reliure)
  outer=2cm,
  top=2.5cm,
  bottom=2.5cm,
  headheight=15pt
]{geometry}
\usepackage{setspace}
\setstretch{1.3}            % interligne lisible (§ 1 du guide)
\sloppy                     % évite les débordements en marge
\usepackage{indentfirst}    % indentation du premier paragraphe après un titre
\setlength{\parindent}{1.2em}
\setlength{\parskip}{0.3em}

% ---------- En-tête / pied de page ------------------------------------
\usepackage{fancyhdr}
\pagestyle{fancy}
\fancyhf{}
\renewcommand{\headrulewidth}{0.4pt}
\renewcommand{\footrulewidth}{0pt}
\fancyhead[L]{\small\nouppercase{\leftmark}}
\fancyhead[R]{\small\thepage}
\fancyfoot[C]{}

% ---------- Couleur et liens ------------------------------------------
\usepackage{xcolor}
\definecolor{supcomblue}{HTML}{1F3A5F}
\definecolor{accent}{HTML}{E8633A}    % orange labConnect
\definecolor{codebg}{HTML}{FAF7F2}
\definecolor{linkblue}{HTML}{1F4E89}

\usepackage[
  colorlinks=true,
  linkcolor=linkblue,
  citecolor=linkblue,
  urlcolor=linkblue,
  pdftitle={Rapport de PFE -- labConnect},
  pdfauthor={Omar Ouatas},
  pdfsubject={Plateforme digitale de gestion de laboratoire d'analyses médicales}
]{hyperref}

% ---------- Titres de chapitres / sections ----------------------------
\usepackage{titlesec}
\titleformat{\chapter}[display]
  {\normalfont\huge\bfseries\color{supcomblue}}
  {\chaptertitlename\ \thechapter}{20pt}{\Huge}
\titlespacing*{\chapter}{0pt}{-30pt}{30pt}

\titleformat{\section}
  {\normalfont\Large\bfseries\color{supcomblue}}
  {\thesection}{1em}{}

\titleformat{\subsection}
  {\normalfont\large\bfseries}
  {\thesubsection}{1em}{}

% ---------- Graphiques, tableaux, code --------------------------------
\usepackage{graphicx}
\graphicspath{{images/}}
\usepackage{caption}
\captionsetup{font=small, labelfont=bf, labelsep=period}
\usepackage{subcaption}
\usepackage{array}
\usepackage{tabularx}
\usepackage{booktabs}
\usepackage{longtable}
\usepackage{multirow}

\usepackage{listings}
\lstdefinestyle{labstyle}{
  backgroundcolor=\color{codebg},
  basicstyle=\ttfamily\footnotesize,
  breaklines=true,
  showstringspaces=false,
  numbers=left,
  numberstyle=\tiny\color{gray},
  frame=single,
  framerule=0pt,
  rulecolor=\color{codebg},
  xleftmargin=2em,
  framexleftmargin=2em,
  framextopmargin=0.5em,
  framexbottommargin=0.5em,
  keywordstyle=\color{supcomblue}\bfseries,
  commentstyle=\color{gray}\itshape,
  stringstyle=\color{accent}
}
\lstset{style=labstyle}

% ---------- Commandes d'aide ------------------------------------------
% Une figure « à compléter » : encadré numéroté avec une légende et un
% espace réservé. Pour insérer une vraie capture, il suffit de remplacer
% \placeholderFigure par \includegraphics dans le code source.
\newcommand{\placeholderFigure}[3][0.75]{%
  \begin{figure}[ht]
    \centering
    \fbox{\begin{minipage}[c][6.5cm][c]{#1\linewidth}
      \centering
      \textit{\small À insérer : #3}
    \end{minipage}}
    \caption{#2}
    \label{fig:#3}
  \end{figure}%
}

% Pour insérer une capture une fois disponible : remplacer le bloc
% \placeholderFigure par :
% \begin{figure}[ht] \centering
%   \includegraphics[width=0.8\linewidth]{nom-fichier}
%   \caption{...} \label{fig:...}
% \end{figure}

% =====================================================================
%                          PAGE DE GARDE
% =====================================================================

\begin{document}

\begin{titlepage}
  \begin{center}
    % ---- Bandeau supérieur : logo SUP'COM ----
    % À INSÉRER : logo officiel SUP'COM (haut de page)
    \fbox{\begin{minipage}[c][2.5cm][c]{0.4\linewidth}
      \centering\textit{\small Logo SUP'COM}
    \end{minipage}}

    \vspace{0.5cm}
    {\large\bfseries École Supérieure des Communications de Tunis}\\[0.4em]
    {\normalsize Cycle des études d'ingénieur en Télécommunications}\\[1.5em]

    {\large Option : \underline{\hspace{6cm}}}\\
    \vspace{2.5cm}

    % ---- Titre principal ----
    {\Huge\bfseries\color{supcomblue} Rapport de Projet}\\[0.4em]
    {\Huge\bfseries\color{supcomblue} de Fin d'Études}\\
    \vspace{2cm}

    % ---- Thème ----
    {\large\bfseries Thème :}\\[0.6em]
    \begin{minipage}{0.85\linewidth}
      \centering\itshape\large
      \textbf{labConnect} — Conception et réalisation d'une plateforme digitale
      de gestion de laboratoire d'analyses médicales adaptée au contexte mauritanien
    \end{minipage}
    \vspace{2cm}

    % ---- Auteur et encadrants ----
    \begin{minipage}[t]{0.45\linewidth}
      \centering
      \textbf{Réalisé par :}\\[0.4em]
      \underline{\hspace{5cm}}\\
      % Ex : M.\ Omar Ouatas
    \end{minipage}%
    \begin{minipage}[t]{0.45\linewidth}
      \centering
      \textbf{Encadrant(s) :}\\[0.4em]
      \underline{\hspace{5cm}}\\
      \underline{\hspace{5cm}}\\
      % Ex : Mme/M. Prénom Nom — SUP'COM
      %      Mme/M. Prénom Nom — Établissement d'accueil
    \end{minipage}
    \vspace{1.5cm}

    \textit{Travail proposé et réalisé en collaboration avec}\\[0.4em]
    \textbf{\underline{\hspace{7cm}}}\\
    % Sigle et nom de l'établissement d'accueil
    \vspace{1.5cm}

    \textbf{Année universitaire : \underline{\hspace{3cm}}}\\

    \vfill
    \rule{\linewidth}{0.5pt}\\[0.4em]
    {\footnotesize
      École Supérieure des Communications de Tunis \textbullet{} SUP'COM\\
      2083 Cité Technologique des Communications, Elghazala, Ariana, Tunisie\\
      Tél. +216 70 240 900 \textbullet{} \href{http://www.supcom.tn}{www.supcom.tn}
    }
  \end{center}
\end{titlepage}

% --- Page blanche après la garde (typo standard) ---
\thispagestyle{empty}\null\newpage

% =====================================================================
%                          DÉDICACES (facultatif)
% =====================================================================
\thispagestyle{empty}
\vspace*{3cm}
\begin{flushright}
  \itshape
  % Quelques lignes personnelles. À compléter par vos soins ; un exemple
  % volontairement sobre, à retoucher selon votre sensibilité.
  À ma famille, qui m'a transmis le goût d'apprendre.\\[0.4em]
  À mes camarades de promotion,\\
  pour les nuits blanches et les éclats de rire.\\[0.4em]
  À celles et ceux qui, en Mauritanie,\\
  travaillent chaque jour à rendre les soins plus accessibles.
\end{flushright}
\newpage

% =====================================================================
%                          RÉSUMÉ ET MOTS-CLÉS
% =====================================================================
\chapter*{Résumé}
\addcontentsline{toc}{chapter}{Résumé}

% Résumé en français — calibré pour ~250 mots. À retoucher pour refléter
% vos propres choix de présentation.
Ce projet de fin d'études porte sur la conception et la réalisation de
\textbf{labConnect}, une plateforme digitale de gestion de laboratoire
d'analyses médicales pensée pour répondre aux besoins du contexte
mauritanien. Le travail part d'un constat simple : la plupart des
solutions de gestion de laboratoire (LIS) disponibles sur le marché sont
coûteuses, conçues pour des structures hospitalières occidentales et ne
prennent pas en compte les spécificités locales — couverture par la
\textit{Caisse Nationale d'Assurance Maladie} (CNAM), faible
informatisation des patients, importance du circuit \emph{walk-in} au
comptoir.

La plateforme livrée se compose de trois briques complémentaires :
une \textbf{API backend} en Django REST, une \textbf{application web}
pour le personnel du laboratoire (chef de labo, biologiste, technicien,
infirmier, secrétaire), et une \textbf{application mobile} pour les
patients et le livreur. Le circuit complet d'une analyse y est
implémenté de bout en bout — création du dossier au comptoir, prélèvement,
saisie du résultat par le technicien, validation par le biologiste avec
possibilité d'édition, et facturation tenant compte du taux de couverture
CNAM du patient. Trois modules transversaux complètent l'ensemble :
un \emph{questionnaire pré-test} qui suit le contexte clinique, un
\emph{tableau de bord statistique} pour le pilotage du laboratoire, et
une \emph{gestion d'inventaire} des réactifs et consommables.

L'architecture retenue privilégie la séparation des préoccupations et
l'auditabilité : chaque opération critique (validation, mouvement de
stock, facturation) laisse une trace horodatée et signée. Le rapport
détaille la démarche, les choix techniques et les résultats obtenus.

\vspace{1em}
\noindent\textbf{Mots-clés :} laboratoire d'analyses médicales, LIS,
Django, React, React Native, CNAM, Mauritanie, multi-rôle, JWT, RBAC,
facturation, gestion d'inventaire.

% --- Abstract anglais ---
\vspace{2em}
\noindent\textbf{\Large Abstract}
\vspace{0.5em}

This final-year engineering project covers the design and implementation
of \textbf{labConnect}, a digital management platform for medical
analysis laboratories, tailored to the Mauritanian context. Most
commercial Laboratory Information Systems (LIS) are expensive and
designed for Western hospital settings; they ignore local specifics such
as the \emph{Caisse Nationale d'Assurance Maladie} (CNAM) coverage rules,
the limited digital footprint of patients, and the importance of
walk-in counter workflows.

The delivered platform combines a Django REST backend, a React web
application for laboratory staff, and a React Native mobile application
for patients and delivery riders. The full lifecycle of an analysis is
implemented end-to-end — counter registration, sample collection,
technician result entry, biologist validation with editing rights, and
invoice issuance accounting for CNAM coverage. Three cross-cutting
modules round out the system: a pre-test clinical questionnaire,
an aggregated statistics dashboard, and a reagent / consumable inventory
manager.

The architecture emphasizes separation of concerns and auditability:
every critical action leaves a timestamped, signed trail.

\vspace{1em}
\noindent\textbf{Keywords:} medical laboratory, LIS, Django, React,
React Native, CNAM, Mauritania, multi-role, JWT, RBAC, billing,
inventory management.

\newpage

% =====================================================================
%                          AVANT-PROPOS
% =====================================================================
\chapter*{Avant-propos}
\addcontentsline{toc}{chapter}{Avant-propos}

% Cadre du projet — à personnaliser
Le présent rapport rend compte du projet de fin d'études que j'ai mené
au cours de l'année universitaire \underline{\hspace{2cm}}, dans le cadre
du cycle des études d'ingénieur en Télécommunications de l'\textbf{École
Supérieure des Communications de Tunis} (SUP'COM). Le projet a été
proposé et accompagné par \underline{\hspace{6cm}}\footnote{Sigle et
nom de l'établissement d'accueil — à compléter.}.

L'idée de \textbf{labConnect} est née d'un constat fait sur le terrain
en Mauritanie : la gestion d'un laboratoire d'analyses médicales y
repose encore très largement sur le papier ou sur des fichiers tableurs
locaux, alors même que les volumes traités et la complexité des
échanges avec la \emph{Caisse Nationale d'Assurance Maladie} ne cessent
de croître. Plutôt que d'essayer d'adapter un produit existant, j'ai
choisi de concevoir une plateforme à partir des contraintes locales,
avec un objectif clair : qu'elle puisse être déployée dans un petit ou
moyen laboratoire sans investissement matériel important, et qu'elle
respecte le découpage métier réel des équipes.

\section*{Remerciements}
% À adapter selon vos encadrants et collègues
Je tiens à exprimer ma sincère gratitude à \textbf{M./Mme
\underline{\hspace{4cm}}}, mon encadrant pédagogique à SUP'COM, pour la
qualité de son suivi, la pertinence de ses remarques et la confiance
qu'il/elle m'a accordée tout au long du projet.

Je remercie également \textbf{M./Mme \underline{\hspace{4cm}}}, mon
encadrant côté \underline{\hspace{3cm}}, pour son accompagnement
quotidien, sa disponibilité et les nombreuses discussions qui m'ont
permis de mieux comprendre les enjeux du métier.

Mes remerciements vont aussi aux membres du jury qui ont accepté
d'évaluer ce travail, ainsi qu'à l'ensemble du corps enseignant de
SUP'COM pour la solidité de la formation reçue.

Enfin, j'adresse une pensée toute particulière à \underline{\hspace{4cm}},
qui m'a fait découvrir le quotidien réel d'un laboratoire d'analyses
médicales en Mauritanie et sans qui ce projet aurait été beaucoup plus
théorique.

\newpage

% =====================================================================
%                          TABLE DES MATIÈRES
% =====================================================================
\renewcommand{\contentsname}{Sommaire}
\tableofcontents
\newpage

\listoffigures
\addcontentsline{toc}{chapter}{Liste des figures}
\newpage

\listoftables
\addcontentsline{toc}{chapter}{Liste des tableaux}
\newpage

% =====================================================================
%                        LISTE DES ABRÉVIATIONS
% =====================================================================
\chapter*{Liste des abréviations}
\addcontentsline{toc}{chapter}{Liste des abréviations}

\begin{tabularx}{\linewidth}{l X}
  \toprule
  \textbf{Sigle} & \textbf{Signification} \\
  \midrule
  API     & Application Programming Interface \\
  CNAM    & Caisse Nationale d'Assurance Maladie (Mauritanie) \\
  CRUD    & Create, Read, Update, Delete \\
  CSS     & Cascading Style Sheets \\
  DRF     & Django REST Framework \\
  HTTP    & HyperText Transfer Protocol \\
  JSON    & JavaScript Object Notation \\
  JWT     & JSON Web Token \\
  LIS     & Laboratory Information System \\
  MRU     & Ouguiya mauritanienne (unité monétaire) \\
  OTP     & One-Time Password \\
  PFE     & Projet de Fin d'Études \\
  PHI     & Protected Health Information \\
  RBAC    & Role-Based Access Control \\
  REST    & Representational State Transfer \\
  RN      & React Native \\
  SDK     & Software Development Kit \\
  SQL     & Structured Query Language \\
  TAT     & Turn-Around Time (délai de rendu d'un résultat) \\
  UI / UX & User Interface / User Experience \\
  URL     & Uniform Resource Locator \\
  UUID    & Universally Unique Identifier \\
  2FA     & Two-Factor Authentication \\
  \bottomrule
\end{tabularx}
\newpage

% =====================================================================
%                       INTRODUCTION GÉNÉRALE
% =====================================================================
\chapter*{Introduction générale}
\addcontentsline{toc}{chapter}{Introduction générale}
\markboth{INTRODUCTION GÉNÉRALE}{}

\section*{Contexte}
Le secteur de la santé connaît, partout dans le monde, une transformation
numérique rapide. Dossier patient informatisé, prescription électronique,
téléconsultation, intelligence artificielle d'aide au diagnostic : les
briques se multiplient et se professionnalisent. Au cœur de ce mouvement,
les laboratoires d'analyses médicales occupent une place particulière —
ils sont à la fois \emph{prestataires} (ils répondent à une demande
émanant d'un médecin ou d'un patient) et \emph{producteurs de données
cliniques} qui irrigueront ensuite la décision thérapeutique.

En Mauritanie, le tableau est contrasté. Le pays s'est doté d'une
\textit{Caisse Nationale d'Assurance Maladie} (CNAM) qui couvre une
partie significative des frais d'analyses pour les assurés, ce qui a
transformé la fréquentation des laboratoires : un patient qui aurait
hésité à dépenser plusieurs milliers d'ouguiyas pour un bilan vient
désormais plus facilement, sachant qu'une partie sera prise en charge.
La conséquence directe est une pression accrue sur les laboratoires de
petite et moyenne taille, dont les outils de gestion sont restés
artisanaux : registres papier, fichiers Excel séparés par poste, échanges
de résultats par WhatsApp. À mesure que les volumes augmentent, ce
système improvisé devient un facteur de risque — erreurs de saisie,
résultats égarés, facturation CNAM imprécise.

\section*{Problématique}
Plutôt que d'importer une solution conçue pour un autre marché — et de
forcer le laboratoire à s'adapter au logiciel — la question initiale de
ce projet a été la suivante : \emph{est-il possible de concevoir une
plateforme de gestion de laboratoire qui parte des contraintes réelles
du contexte mauritanien, qui respecte le découpage métier des équipes
sur place, et qui reste accessible à une petite structure ?}

Cette question en cache plusieurs autres :

\begin{itemize}
  \item Comment modéliser le cycle de vie d'une analyse, depuis l'arrivée
        du patient jusqu'à la publication du résultat validé, en
        impliquant correctement les cinq rôles métier (secrétaire,
        infirmier·e, technicien·ne, biologiste, chef de laboratoire) ?
  \item Comment intégrer nativement la couverture CNAM dans le calcul
        de la facture, sans forcer un re-paramétrage manuel à chaque
        passage ?
  \item Comment garantir la traçabilité des résultats — qui les a
        saisis, qui les a validés, qui les a éventuellement corrigés —
        sans alourdir l'expérience utilisateur ?
  \item Comment permettre au patient, lui aussi, d'utiliser la
        plateforme à son rythme, via une application mobile, sans rendre
        cet usage obligatoire ?
\end{itemize}

\section*{Présentation du sujet}
La réponse construite dans ce projet est une plateforme à trois étages,
nommée \textbf{labConnect}, qui couvre :

\begin{enumerate}
  \item Une \textbf{API backend} en Django + Django REST Framework,
        servant de source de vérité pour les données métier et exposant
        une cinquantaine d'endpoints sécurisés par JWT.
  \item Une \textbf{application web} pour le personnel du laboratoire,
        développée en React 19 + Vite, qui adapte ses écrans aux
        permissions de l'utilisateur connecté (un même biologiste qui
        est aussi technicien voit l'union des deux périmètres).
  \item Une \textbf{application mobile} pour le patient et le livreur,
        développée en React Native (Expo SDK 54), qui permet de
        rechercher un laboratoire, prendre rendez-vous, consulter ses
        résultats validés, et — côté livreur — assurer la collecte
        d'échantillons à domicile.
\end{enumerate}

\section*{Plan du rapport}
Le rapport est organisé en quatre chapitres.

Le \textbf{chapitre 1} dresse le cadre du projet : il présente le
contexte mauritanien, analyse les solutions LIS existantes, formalise
les besoins fonctionnels et non fonctionnels exprimés par les équipes
des laboratoires, et trace les grandes lignes de l'architecture cible.

Le \textbf{chapitre 2} décrit la conception détaillée : modèle de
données, machine à états du cycle de vie d'une analyse, gestion
multi-rôle par groupes Django, stratégie d'authentification à deux
facteurs et règles métier transverses (CNAM, traçabilité, audit).

Le \textbf{chapitre 3} se concentre sur la réalisation : la stack
technique retenue, les choix d'implémentation, les composants
réutilisables côté frontend, la gestion des migrations de base de
données, et la mise en \oe{}uvre des fonctionnalités phares (circuit
analyse complet, facturation CNAM, questionnaire pré-test, statistiques,
inventaire).

Le \textbf{chapitre 4} présente la validation et les résultats : tests
fonctionnels, scénarios de bout en bout, captures de l'application en
contexte d'usage réel, retours utilisateurs et perspectives.

Le rapport se clôt par une conclusion générale qui revient sur les
contributions du projet et ouvre des perspectives concrètes
d'évolution.

\newpage

% =====================================================================
%                          CHAPITRE 1 — CADRE
% =====================================================================
\chapter{Cadre du projet et étude préalable}

% Introduction de chapitre — chaque chapitre s'ouvre par un court
% paragraphe qui explique son rôle dans la progression du rapport.
\noindent
Avant de présenter la solution, il faut comprendre le contexte qui l'a
fait émerger. Ce chapitre dresse le décor : il revient sur la
situation du secteur de la santé en Mauritanie, examine les outils
existants côté laboratoire, formalise les besoins exprimés par les
équipes que j'ai pu rencontrer, et trace l'architecture macroscopique
cible. C'est le socle sur lequel la conception détaillée du
chapitre~\ref{ch:conception} viendra s'appuyer.

\section{Contexte du secteur de la santé en Mauritanie}

% --- Sous-section : portrait général ---
\subsection{Une couverture santé qui se structure}
La République Islamique de Mauritanie s'est dotée, à partir des années
2000, d'un dispositif de couverture santé articulé autour de la
\textbf{Caisse Nationale d'Assurance Maladie} (CNAM). La couverture
proposée pour les analyses biomédicales varie selon le statut de
l'assuré, mais s'établit usuellement autour de 80~\% des frais
engagés~\cite{cnam2024}. Le reste à charge — la part patient — est
réglé directement au comptoir du laboratoire.

Cette répartition simple a une conséquence importante pour les
laboratoires : ils ne facturent jamais un seul montant à un seul
destinataire. Pour chaque analyse, ils doivent identifier la fraction
prise en charge (à transmettre à la CNAM pour remboursement) et la
fraction restant due (à encaisser auprès du patient). Mal géré, ce
double flux est une source importante de pertes financières : un
laboratoire qui se trompe d'un point de pourcentage sur quelques milliers
d'analyses par an perd vite plusieurs centaines de milliers d'ouguiyas.

% --- À INSÉRER : carte ou photo illustrative du contexte ---
\placeholderFigure{Carte de la couverture santé en Mauritanie}%
{contexte-mauritanie-carte}{capture illustrative — carte de Nouakchott
avec localisation des principaux laboratoires partenaires (par
exemple via OpenStreetMap)}

\subsection{Un parc de laboratoires hétérogène}
Au-delà du CHU et des grandes structures hospitalières publiques, le
paysage est dominé par des laboratoires privés de petite et moyenne
taille (de 3 à 15 employés). Leur équipement informatique se limite
souvent à un ordinateur de bureau partagé entre la secrétaire et le
chef de laboratoire. Quand un logiciel existe, c'est en général une
application monoposte développée localement ou un fichier Excel
sophistiqué.

% --- À INSÉRER : photo prise dans un laboratoire mauritanien (avec
%     autorisation) montrant un poste de travail typique : registre
%     papier + ordinateur. ---
\placeholderFigure{Poste de travail d'une secrétaire dans un laboratoire de Nouakchott}%
{contexte-poste-secretaire}{photo de terrain — poste de travail
mêlant registre papier et ordinateur, avec accord du laboratoire}

\section{État de l'art : les LIS existants}

\subsection{Définition et fonctions d'un LIS}
Un \emph{Laboratory Information System} (LIS) est un logiciel qui
prend en charge l'ensemble du cycle de vie d'une analyse, depuis
l'enregistrement de la demande jusqu'à la diffusion du résultat
validé~\cite{tate2014lis}. Les briques fonctionnelles typiques sont :

\begin{itemize}
  \item Gestion du dossier patient et des prescriptions.
  \item Suivi des échantillons (étiquetage, réception, rejet).
  \item Saisie et validation des résultats.
  \item Connexion aux automates analytiques (interface ASTM ou HL7).
  \item Facturation et reporting financier.
\end{itemize}

\subsection{Solutions du marché et limites}
Les LIS commerciaux établis (Cerner CoPath, Sunquest, LabWare, etc.)
sont des produits éprouvés, riches en fonctionnalités, mais qui posent
trois problèmes dans le contexte qui nous occupe :

\begin{enumerate}
  \item \textbf{Coût d'acquisition et de maintenance.} Les licences se
        chiffrent en dizaines de milliers d'euros, sans compter le
        matériel serveur requis. Hors de portée d'un laboratoire
        mauritanien type.
  \item \textbf{Inadéquation culturelle.} Les workflows sont calés sur
        des hôpitaux occidentaux, avec un médecin prescripteur
        systématique. Or, en Mauritanie, une part importante des
        patients arrive au comptoir avec une demande directe, sans
        ordonnance.
  \item \textbf{Absence de support CNAM.} La gestion des assureurs
        proposée par défaut suppose des contrats de tiers payant
        complexes ; rien n'est prévu pour la mécanique simple
        \emph{taux de couverture × prix unitaire} qui correspond à la
        CNAM.
\end{enumerate}

\subsection{Solutions open source}
Du côté libre, OpenELIS, BlueRedox et quelques projets dérivés
d'OpenMRS proposent des LIS adaptables. Ces solutions sont
intéressantes mais souffrent d'une dette UX visible (interfaces datées,
parcours utilisateurs lourds) et nécessitent un effort d'intégration
non négligeable. Le coût total d'adoption — formation, paramétrage,
adaptation — reste élevé pour une petite structure.

\section{Analyse des besoins}

% À INSÉRER éventuellement : tableau-photo de prise de notes lors des
% entretiens utilisateurs réalisés sur le terrain.

\subsection{Méthodologie}
Pour cerner les besoins réels, j'ai mené une série de courts entretiens
semi-directifs avec des professionnels exerçant en Mauritanie :
deux chefs de laboratoire, un biologiste, deux techniciennes et une
secrétaire médicale. Les entretiens, d'une durée de 30 à 45 minutes,
portaient sur cinq axes :

\begin{enumerate}
  \item Déroulement type d'une journée.
  \item Outils utilisés et points de friction.
  \item Gestion de la CNAM.
  \item Pratiques d'archivage et de partage des résultats.
  \item Souhaits exprimés pour un outil idéal.
\end{enumerate}

\subsection{Besoins fonctionnels}
La synthèse des entretiens permet de regrouper les besoins fonctionnels
en six familles, présentées dans le tableau~\ref{tab:besoins-fonctionnels}.

\begin{table}[ht]
  \centering
  \caption{Synthèse des besoins fonctionnels exprimés par les équipes interrogées}
  \label{tab:besoins-fonctionnels}
  \begin{tabularx}{\linewidth}{l X}
    \toprule
    \textbf{Famille} & \textbf{Besoin formulé} \\
    \midrule
    Patient        & Enregistrement rapide au comptoir (\emph{walk-in}) avec création automatique du dossier ; consultation des résultats par le patient depuis son téléphone. \\
    Analyses       & Catalogue de tests configurable, prix figés à la commande, questionnaire pré-test paramétrable par test. \\
    Workflow       & Saisie du résultat par le technicien ; validation par le biologiste avec possibilité de corriger la valeur saisie tout en gardant trace de l'original. \\
    Facturation    & Calcul automatique du split CNAM / patient à la création de l'ordre ; édition d'une facture détaillée imprimable. \\
    Pilotage       & Tableau de bord d'activité (volumes, statuts, délais) avec sélecteur de fenêtre temporelle ; alertes en cas de stock bas. \\
    Inventaire     & Suivi des réactifs et consommables avec mouvements tracés (livraison, consommation, ajustement). \\
    \bottomrule
  \end{tabularx}
\end{table}

\subsection{Besoins non fonctionnels}
Quatre besoins non fonctionnels structurants ont été identifiés :

\begin{itemize}
  \item \textbf{Multi-rôle.} Dans les petits laboratoires, une même
        personne porte souvent plusieurs casquettes : un technicien qui
        accueille aussi les patients quand la secrétaire est absente,
        un biologiste qui saisit lui-même certains résultats. La
        plateforme doit gérer nativement ce cumul de rôles sans imposer
        de bascule manuelle de contexte.
  \item \textbf{Auditabilité.} Toute opération critique (saisie de
        résultat, validation, mouvement de stock, facturation) doit
        laisser une trace horodatée et associée à un utilisateur
        identifié.
  \item \textbf{Sécurité.} Les données patient relèvent du \emph{Protected
        Health Information} (PHI). L'authentification doit reposer sur
        des standards éprouvés (JWT signés, hash des mots de passe,
        authentification à deux facteurs côté staff).
  \item \textbf{Faible empreinte.} La plateforme doit fonctionner sur
        un serveur modeste (4 Go de RAM suffisent pour un petit
        laboratoire) et sur des smartphones d'entrée de gamme côté
        mobile.
\end{itemize}

\section{Choix d'une architecture cible}

% À INSÉRER : schéma d'architecture macroscopique — trois clients
% (web staff, mobile patient, mobile livreur) reliés au backend Django
% qui parle à PostgreSQL ; sur le côté, intégration optionnelle avec
% Firebase pour l'authentification sociale et avec OpenStreetMap pour
% la cartographie.
\placeholderFigure{Architecture macroscopique de la plateforme labConnect}%
{architecture-macro}{schéma d'architecture — trois clients (web staff,
mobile patient, mobile livreur) connectés au backend Django, base
PostgreSQL, intégrations Firebase et OpenStreetMap}

\subsection{Une plateforme à trois étages}
Trois choix architecturaux fondamentaux ont structuré la suite du
projet :

\begin{enumerate}
  \item \textbf{Un backend unique servant toutes les UIs.} Le backend
        Django expose une API REST unique consommée à la fois par le
        frontend web et par l'application mobile. Cela évite la
        duplication des règles métier — la logique de validation d'un
        résultat est définie une seule fois.
  \item \textbf{Une séparation stricte des rôles entre apps.} Le web
        est pour le personnel du laboratoire ; le mobile pour le
        patient et le livreur. Aucun écran « staff » n'existe sur
        mobile, et inversement. Cela permet d'optimiser l'UX pour
        chaque audience.
  \item \textbf{Une approche multi-laboratoire dès le départ.} Le
        modèle de données est multi-tenant : un chef de laboratoire
        peut gérer plusieurs sites, et un utilisateur peut être
        rattaché à plusieurs laboratoires. L'isolation se fait par le
        champ \texttt{laboratory\_id} sur chaque entité métier.
\end{enumerate}

\subsection{Stack technique retenue}
Le tableau~\ref{tab:stack} récapitule la stack technique retenue, avec
le rôle de chaque brique et la justification du choix.

\begin{table}[ht]
  \centering
  \caption{Stack technique de la plateforme labConnect}
  \label{tab:stack}
  \small
  \begin{tabularx}{\linewidth}{l l X}
    \toprule
    \textbf{Couche} & \textbf{Technologie} & \textbf{Justification} \\
    \midrule
    Backend        & Django 5 + DRF              & Cadriciel mature, fort sur la sécurité et l'administration. ORM solide. \\
    Base de données & PostgreSQL + PostGIS       & Robustesse transactionnelle ; PostGIS pour géolocalisation des labos. \\
    Authentification & JWT (SimpleJWT)           & Stateless, parfait pour servir un web et un mobile depuis la même API. \\
    Frontend web   & React 19 + Vite             & Performances, écosystème, hot reload rapide en développement. \\
    Frontend mobile & React Native + Expo SDK 54 & Capitalisation des compétences React, déploiement OTA possible. \\
    Cartographie   & Leaflet + OpenStreetMap     & Pas de clé API à payer, données ouvertes. \\
    Identité sociale & Firebase Auth             & Délégation propre du login Google / Email pour le patient mobile. \\
    \bottomrule
  \end{tabularx}
\end{table}

\section*{Conclusion du chapitre}
Ce chapitre a posé le décor. Le secteur de la santé mauritanien est en
pleine structuration, mais ses laboratoires de petite taille manquent
d'outils adaptés. Les LIS commerciaux sont trop lourds et trop chers,
les LIS open source existent mais n'intègrent pas la mécanique CNAM
native. Les entretiens utilisateurs ont permis de formaliser six
familles de besoins fonctionnels et quatre exigences non fonctionnelles
qui guident la conception. Le chapitre suivant entre dans le détail de
cette conception.

\newpage

% =====================================================================
%                  CHAPITRE 2 — CONCEPTION DÉTAILLÉE
% =====================================================================
\chapter{Conception détaillée}\label{ch:conception}

\noindent
Ce chapitre traduit les besoins en modèles. Il décrit le modèle de
données, la machine à états du cycle de vie d'une analyse, la gestion
multi-rôle, la stratégie d'authentification, et les règles métier
transverses qui s'appliqueront ensuite à toutes les fonctionnalités.

\section{Modèle de données}

\subsection{Diagramme d'entité-relation}
Le schéma de la base de données s'articule autour de cinq entités
principales : \texttt{User}, \texttt{Laboratory}, \texttt{Appointment},
\texttt{Sample} et \texttt{TestOrder} — chacune complétée par des
profils ou des entités secondaires (résultats, mouvements
d'inventaire, items du catalogue).

% À INSÉRER : diagramme UML / E-R généré par exemple avec dbdiagram.io
% ou draw.io. Cible : une page A4 paysage avec toutes les tables
% principales et les relations clés (1-N, M-N).
\placeholderFigure[0.95]{Diagramme entité-relation principal de labConnect}%
{er-diagram}{diagramme E-R — généré depuis dbdiagram.io ou draw.io,
montrant Laboratory, User, PatientProfile, StaffProfile, Appointment,
Sample, TestCatalogEntry, TestOrder, TestResult, InventoryItem,
InventoryMovement}

\subsection{Choix de modélisation marquants}

\paragraph{Profils séparés.} La table \texttt{User} reste minimale
(identité, contact, sécurité). Les attributs métier vivent dans des
\emph{profils} satellites : \texttt{PatientProfile} (date de naissance,
groupe sanguin, numéro CNAM\ldots) et \texttt{StaffProfile} (un par
couple utilisateur-laboratoire, ce qui permet à un même utilisateur
d'être rattaché à plusieurs sites). Cette séparation évite d'embarquer
les champs PHI dans la table d'authentification.

\paragraph{UUID exposés, IDs internes cachés.} Toutes les entités
héritent d'un \texttt{BaseModel} qui leur ajoute un champ
\texttt{uuid} (UUIDv4) et des timestamps. C'est l'UUID qui circule
dans l'API ; le \texttt{BigAutoField} interne reste invisible. Cela
évite les attaques par énumération et facilite une éventuelle
migration de base.

\paragraph{Soft delete par défaut.} \texttt{BaseModel} ajoute aussi un
champ \texttt{deleted\_at}. Un \emph{manager} \texttt{active} exclut
automatiquement les lignes supprimées. Les opérations destructives
deviennent réversibles, ce qui est précieux dans un contexte médical
où l'effacement définitif est souvent interdit.

\paragraph{Split CNAM figé.} Sur chaque ligne \texttt{TestOrder}, le
prix unitaire et le split \texttt{cnam\_covered\_mru} /
\texttt{patient\_due\_mru} sont \emph{figés au moment de la création}.
Si la CNAM change ses taux dans six mois, les factures historiques ne
sont pas mutées rétroactivement. Cette décision, anodine sur le
papier, est en réalité fondamentale pour la conformité comptable.

\section{Machine à états du cycle de vie d'une analyse}

% À INSÉRER : diagramme d'états-transitions montrant les 5 statuts
% d'un TestOrder : pending → in_progress → completed → validated, avec
% transition latérale possible vers rejected.
\placeholderFigure{Machine à états d'un TestOrder}%
{statemachine-order}{diagramme d'états — pending, in\_progress,
completed, validated, rejected, avec les transitions et le rôle
autorisé pour chaque transition}

\subsection{Les cinq statuts}
Un \texttt{TestOrder} se déplace entre cinq statuts canoniques :

\begin{enumerate}
  \item \textbf{pending} — créé, en attente du prélèvement.
  \item \textbf{in\_progress} — l'échantillon a été réceptionné par
        un technicien, la mesure peut être lancée.
  \item \textbf{completed} — le technicien a saisi un résultat ; il
        attend la validation du biologiste.
  \item \textbf{validated} — le biologiste a validé. Le résultat est
        publié au patient.
  \item \textbf{rejected} — l'échantillon a été rejeté (hémolyse,
        identification erronée, etc.).
\end{enumerate}

\subsection{Acteurs autorisés par transition}
Chaque transition est gardée par une vérification de rôle au niveau
de la vue. Le tableau~\ref{tab:transitions} récapitule les
permissions.

\begin{table}[ht]
  \centering
  \caption{Acteurs autorisés à déclencher chaque transition}
  \label{tab:transitions}
  \begin{tabularx}{\linewidth}{l l X}
    \toprule
    \textbf{Transition} & \textbf{Méthode HTTP} & \textbf{Rôles autorisés} \\
    \midrule
    pending → in\_progress       & PATCH \texttt{/lab/samples/.../receive/} & Technicien, biologiste, chef \\
    in\_progress → completed     & POST \texttt{/lab/orders/.../result/}    & Technicien, biologiste, chef \\
    completed → validated        & PATCH \texttt{/lab/orders/.../result/}   & Biologiste, chef \\
    * → rejected                 & PATCH \texttt{/lab/samples/.../reject/}  & Technicien, biologiste, chef \\
    \bottomrule
  \end{tabularx}
\end{table}

\section{Gestion multi-rôle}

\subsection{Le choix des groupes Django}
Plutôt que d'introduire une couche RBAC maison, je me suis appuyé sur
les groupes natifs de Django. Chaque rôle métier (\texttt{secretary},
\texttt{nurse}, \texttt{technician}, \texttt{biologist},
\texttt{lab\_chief}, \texttt{patient}) correspond à un groupe ; un
utilisateur peut appartenir à plusieurs groupes simultanément.

Cette décision a deux avantages immédiats. D'une part, on hérite
gratuitement de l'interface d'administration Django pour la gestion
des rôles. D'autre part, le contrôle d'accès dans les vues se
résume à une intersection d'ensembles.

\subsection{Union des permissions}
Côté frontend web, j'ai défini une table de \emph{permissions atomiques}
(\texttt{enterResult}, \texttt{validate}, \texttt{createSample},
\texttt{viewFinance}, \texttt{editTests}, \texttt{editStaff},
\texttt{editSettings}) qui est calculée comme l'\textbf{union} des
permissions associées à chaque rôle de l'utilisateur. Concrètement, un
utilisateur biologiste \emph{et} technicien voit aussi bien la file de
saisie que la file de validation, sans avoir à basculer manuellement
d'un mode à l'autre.

\section{Authentification}

% À INSÉRER : diagramme de séquence du login JWT — client envoie
% credentials, backend renvoie access + refresh, client utilise access
% dans Authorization header, refresh quand 401.
\placeholderFigure{Diagramme de séquence du login JWT}%
{seq-jwt}{diagramme de séquence — login JWT avec rafraîchissement
automatique du token expiré côté client}

\subsection{JWT pour staff et patient}
L'API utilise des JSON Web Tokens (JWT) signés via la bibliothèque
\texttt{rest\_framework\_simplejwt}. Deux jetons sont émis :
un \texttt{access} de courte durée (15 minutes) et un \texttt{refresh}
de longue durée (7 jours). Le payload de l'\texttt{access} contient
le \texttt{user\_id}, la liste des rôles et un flag
\texttt{totp\_verified} — ce qui permet aux vues de vérifier les
permissions sans interroger la base à chaque requête.

\subsection{Identité sociale via Firebase}
Pour l'application mobile patient, j'ai intégré Firebase Auth comme
\emph{fournisseur d'identité} (et non comme système de session
persistante). Le flux est le suivant :

\begin{enumerate}
  \item Le patient se connecte via Email/Google côté Firebase et
        obtient un \emph{ID token} signé par Google.
  \item L'application mobile transmet cet ID token au backend.
  \item Le backend vérifie la signature (en téléchargeant les clés
        publiques Google) et, si valide, crée ou récupère le compte
        patient correspondant.
  \item Le backend émet ses propres JWT \texttt{access} / \texttt{refresh},
        qui deviennent la véritable session.
\end{enumerate}

Cette architecture évite de dépendre de Firebase pour la persistance
côté mobile (et l'avertissement \emph{AsyncStorage persistence}
caractéristique) tout en bénéficiant de la qualité de leur SDK
d'authentification sociale.

\subsection{2FA mockée en démo}
En contexte de démonstration et d'évaluation, l'envoi effectif d'un
code OTP par SMS est désactivé : le backend force le flag
\texttt{totp\_verified} à \texttt{true} pour tous les utilisateurs.
Les endpoints TOTP existent toujours, mais ne bloquent rien. La
remise en production réelle se fera en supprimant ce \emph{mock}
dans le module \texttt{accounts/tokens.py}.

\section{Règles métier transverses}

\subsection{Traçabilité de la validation}
Quand un biologiste valide un résultat, il peut aussi le \emph{corriger}
— par exemple parce qu'une relecture critique conduit à requalifier le
flag « normal » en « haut ». Pour ne pas perdre la trace de la saisie
originale du technicien, le modèle \texttt{TestResult} embarque un
champ \texttt{original\_value}. Ce champ reste vide tant que personne
n'a corrigé. Dès qu'un biologiste modifie effectivement une valeur, le
backend snapshote la version précédente sous forme lisible
(« \textit{12.4 mUI/L [normal]} ») et l'écrit dans
\texttt{original\_value}. Si un autre biologiste corrige par la suite,
ce champ n'est pas écrasé — il pointe toujours sur la version
\emph{technicien}.

\subsection{Mouvements de stock atomiques}
Le module inventaire suit une règle stricte : le champ
\texttt{current\_stock} d'un \texttt{InventoryItem} n'est jamais
modifié directement. Toute variation passe par un
\texttt{InventoryMovement} muni d'un \texttt{delta} signé et d'une
raison (\texttt{delivery}, \texttt{consumption}, \texttt{adjustment},
\texttt{expiry}). La création du mouvement et la mise à jour du stock
sont enveloppées dans une transaction PostgreSQL avec
\texttt{SELECT FOR UPDATE} sur la ligne concernée, ce qui évite les
\emph{data races} si deux agents mouvementent le même article
simultanément.

\subsection{Facturation et CNAM}
Le calcul du split CNAM utilise un arrondi à l'entier inférieur côté
CNAM, le reste éventuel étant à la charge du patient. Le code
correspondant est court et tient en cinq lignes Python\,:

\begin{lstlisting}[language=Python,caption={Calcul du split CNAM / patient (analyses/views.py)},label={lst:split-cnam}]
def _split_cnam(price_mru, coverage_pct: int) -> tuple[int, int]:
    p = int(price_mru)
    pct = max(0, min(100, int(coverage_pct or 0)))
    covered = (p * pct) // 100
    return covered, p - covered
\end{lstlisting}

\noindent
Le choix de l'arrondi vers le bas pour la CNAM est délibéré : sur un
éventuel résidu d'arrondi, on charge plutôt le patient (centime près)
que la sécurité sociale, pour ne jamais sur-facturer la CNAM à cause
d'une approximation numérique.

\section*{Conclusion du chapitre}
La conception détaillée s'organise autour de quelques décisions
structurantes : un modèle de données minimal mais expressif (profils
séparés, UUID exposés, soft delete), une machine à états explicite
pour le cycle de vie de l'analyse, une gestion multi-rôle par union
des permissions, et une mécanique d'authentification mixte JWT-natif
plus Firebase social. Le chapitre suivant entre dans la réalisation.

\newpage

% =====================================================================
%                       CHAPITRE 3 — RÉALISATION
% =====================================================================
\chapter{Réalisation}

\noindent
La conception une fois figée, vient l'étape la plus longue : la
réalisation. Ce chapitre revient sur l'organisation du code, les
patterns réutilisés à travers les écrans, et la mise en \oe{}uvre des
quatre fonctionnalités les plus structurantes du projet : le circuit
analyse complet, la facturation CNAM, le questionnaire pré-test, et
les modules transverses statistiques et inventaire.

\section{Organisation du dépôt}

% À INSÉRER : capture du dépôt Git (par exemple un screenshot de
% l'arborescence dans VS Code ou un graphique généré par tree).
\placeholderFigure{Arborescence du dépôt labConnect}%
{repo-tree}{capture VS Code — arborescence du dépôt avec les trois
sous-projets Backend, frontend-web, frontend-client}

Le projet est organisé en trois sous-dépôts coexistant dans un
\emph{monorepo} simple :

\begin{itemize}
  \item \texttt{Backend/} — projet Django et ses cinq applications
        (\texttt{core}, \texttt{accounts}, \texttt{laboratories},
        \texttt{appointments}, \texttt{analyses}, \texttt{inventory}).
  \item \texttt{frontend-web/} — application web Vite + React 19
        pour le personnel.
  \item \texttt{frontend-client/} — application Expo React Native
        pour le patient et le livreur.
\end{itemize}

La documentation centrale (\texttt{API.md}, \texttt{FIREBASE\_SETUP.md})
vit au niveau racine. Ce monorepo facilite la cohérence entre les
contrats d'API et les clients qui les consomment.

\section{Backend Django : patterns récurrents}

\subsection{BaseModel et soft delete}
Le module \texttt{core/models.py} définit un \texttt{BaseModel}
abstrait. Toutes les entités métier en héritent, ce qui leur fournit
gratuitement un UUID, des timestamps, le soft delete et le manager
\texttt{active}. Le bloc tient en une vingtaine de lignes mais
économise des centaines de lignes répétitives à l'échelle de
l'application.

\subsection{Permissions avec scope laboratoire}
Le module \texttt{core/permissions.py} expose une classe
\texttt{IsStaffAnd2FA} et une fonction \texttt{current\_lab\_id(request)}
qui lit l'en-tête HTTP \texttt{X-Lab-Uuid} envoyé par le client. Toutes
les vues métier filtrent leurs querysets par ce \texttt{lab\_id}, ce
qui garantit l'isolation entre laboratoires sans aucune ligne de code
spécifique dans les vues.

\subsection{Sérialiseurs DRF}
J'ai suivi la convention DRF \emph{un sérialiseur lecture, un
sérialiseur écriture} pour les vues complexes : par exemple, le
catalogue de tests dispose d'un \texttt{TestCatalogSerializer} pour
les listes et d'un \texttt{TestCatalogCreateSerializer} pour la
création. Cela évite les champs en \texttt{read\_only}/\texttt{write\_only}
qui rendent les sérialiseurs difficiles à lire.

\section{Frontend web : composants partagés}

\subsection{Modal réutilisable}
Plutôt que de recoder une boîte modale à chaque écran, j'ai centralisé
le composant \texttt{Modal.jsx}. Il expose une API uniforme (titre,
sous-titre, footer, taille \texttt{wide}) et gère seul l'overlay, la
fermeture par clic extérieur et l'ESC. Tous les formulaires importants
de l'application (création de test, walk-in, validation, facture,
mouvement d'inventaire) l'utilisent.

\subsection{Icônes inline}
L'application embarque une bibliothèque d'icônes définies en JSX
directement dans \texttt{icons.jsx}. Aucune dépendance externe, aucune
charge réseau supplémentaire. Le tableau~\ref{tab:icons} montre la
trentaine d'icônes utilisées.

\begin{table}[ht]
  \centering
  \caption{Extrait de la bibliothèque d'icônes inline}
  \label{tab:icons}
  \begin{tabularx}{\linewidth}{l l X}
    \toprule
    \textbf{Nom} & \textbf{Famille} & \textbf{Utilisation principale} \\
    \midrule
    Home      & Navigation & Tableau de bord \\
    Sparkle   & Navigation & Écran Analyses \\
    Beaker    & Navigation & Catalogue de tests \\
    BarChart  & Navigation & Statistiques \\
    Box       & Navigation & Inventaire \\
    Check     & Action     & Validation, soumission \\
    Edit      & Action     & Modification d'un item \\
    Trash     & Action     & Suppression \\
    Lock      & État       & Permission refusée \\
    \bottomrule
  \end{tabularx}
\end{table}

\section{Fonctionnalité phare : le circuit analyse complet}

\subsection{Walk-in au comptoir}
Le besoin remonté par les secrétaires était clair : permettre la
création d'une analyse pour un patient qui se présente \emph{sans
avoir l'application mobile}. La solution implémentée est un endpoint
unique, \texttt{POST /lab/walk-in/}, qui réalise atomiquement quatre
opérations :

\begin{enumerate}
  \item \texttt{find-or-create} du patient par son numéro de téléphone.
  \item Création d'un \texttt{Appointment} \texttt{visit\_type=in\_lab}
        et \texttt{status=confirmed}, programmé pour l'instant présent.
  \item Création d'un \texttt{Sample} avec un code-barres
        auto-généré.
  \item Création d'un \texttt{TestOrder} par test sélectionné, avec
        prix figé.
\end{enumerate}

Si l'une des opérations échoue, la transaction est annulée
intégralement.

% À INSÉRER : capture du modal Walk-in dans l'app web — montrer
% les 3-4 sections (Patient, Tests, Questionnaire, Notes) avec un
% patient en cours de création.
\placeholderFigure{Modal Walk-in dans l'application web}%
{walkin-modal}{capture écran — modal « Nouvelle analyse au comptoir »
avec sections Patient, Tests, CNAM, Questionnaire pré-test, Notes}

\subsection{Saisie technicien}
Lorsque le technicien réceptionne l'échantillon, il bascule
l'\texttt{Appointment} vers \texttt{in\_progress}. La saisie du
résultat se fait via un modal dédié qui collecte la valeur, l'unité,
la plage de référence, le drapeau (normal / bas / haut / critique) et
d'éventuelles notes techniques.

% À INSÉRER : capture du modal de saisie technicien avec un cas
% concret (par exemple TSH 2.3 mUI/L flag normal).
\placeholderFigure{Saisie d'un résultat par le technicien}%
{result-entry}{capture écran — modal de saisie résultat avec valeur,
unité, plage de référence, drapeau et notes}

\subsection{Validation biologiste avec édition}
Le biologiste, quand il valide, peut soit accepter tel quel le résultat
du technicien, soit le \emph{corriger}. La carte du haut du modal
affiche en lecture seule \emph{qui} a saisi le résultat et
\emph{quand}. En dessous, les champs sont éditables. Si une
modification est détectée, un bandeau d'avertissement indique que la
valeur d'origine sera archivée.

% À INSÉRER : capture du modal de validation biologiste, idéalement
% avec une correction en cours (valeur modifiée par rapport à la
% saisie du technicien).
\placeholderFigure{Validation biologiste avec édition}%
{validation-modal}{capture écran — modal de validation avec
correction de valeur et bandeau « valeur d'origine sera archivée »}

\subsection{Facturation CNAM}
Une fois le résultat validé, le bouton « Facture » devient accessible
aux utilisateurs disposant de la permission \texttt{viewFinance}.
Il ouvre un modal qui affiche le détail item par item, avec la part
CNAM et la part patient pour chaque test, et les totaux en bas.
Un bouton « Imprimer » ouvre la vue d'impression native du navigateur.

% À INSÉRER : capture du modal « Facture » montrant un patient
% couvert CNAM 80 % avec deux ou trois tests.
\placeholderFigure{Facture détaillée d'un rendez-vous}%
{invoice-modal}{capture écran — modal facture avec lignes par test,
split CNAM/patient et totaux en bas}

\section{Fonctionnalité phare : questionnaire pré-test}

Le questionnaire est paramétrable au niveau de chaque entrée du
catalogue. Le chef de laboratoire saisit ses questions sous forme
d'une liste textuelle (\emph{une question par ligne}) ; le backend
les stocke en JSON. Au moment du walk-in, les questions associées
aux tests sélectionnés sont affichées et les réponses sont collectées
puis attachées à chaque \texttt{TestOrder}. Pour ne pas être perturbée
par une évolution ultérieure du catalogue, la liste des questions est
snapshotée sur l'ordre au moment de la création.

% À INSÉRER : capture combinée — à gauche, l'édition du questionnaire
% dans le modal Tests ; à droite, l'affichage des réponses dans le
% modal de validation côté biologiste.
\placeholderFigure[0.95]{Cycle du questionnaire pré-test}%
{questionnaire-flow}{capture combinée — édition du questionnaire dans
le catalogue, saisie au walk-in, affichage dans la validation
biologiste}

\section{Module transverse : statistiques}

L'endpoint \texttt{GET /lab/stats/?days=N} agrège l'activité du
laboratoire sur une fenêtre glissante. La réponse contient cinq
sections :

\begin{itemize}
  \item \texttt{kpis} — quatre indicateurs principaux (ordres totaux,
        validés, délai moyen, revenu).
  \item \texttt{by\_status}, \texttt{by\_category}, \texttt{by\_day} —
        trois répartitions visualisables sous forme de barres.
  \item \texttt{by\_technician} — \emph{throughput} individuel.
  \item \texttt{low\_stock\_items} — clin d'\oe{}il vers l'inventaire,
        pour mettre en lumière les ruptures imminentes.
\end{itemize}

Le backend filtre les chiffres financiers selon le rôle : un
technicien voit son activité mais pas le chiffre d'affaires.

% À INSÉRER : capture de l'écran Statistiques avec la période
% sélectionnée à 30 jours, montrant les KPI cards en haut et les bar
% charts en dessous.
\placeholderFigure{Tableau de bord statistique}%
{stats-screen}{capture écran — Stats.jsx avec fenêtre 30 jours, KPIs
en haut, bar charts par jour et par statut, et carte stock bas en bas
si applicable}

\section{Module transverse : gestion d'inventaire}

\subsection{Modèle item + mouvements}
Le module \texttt{inventory} expose deux entités :
\texttt{InventoryItem} (description et stock courant) et
\texttt{InventoryMovement} (variation signée et raison). Le stock
courant n'est jamais modifié directement — toute mise à jour passe
par la méthode de classe \texttt{InventoryMovement.apply()}, qui
réalise atomiquement la création du mouvement et la mise à jour du
stock dans une transaction PostgreSQL.

\subsection{Quatre types de mouvements}
Quatre raisons sont possibles, chacune avec ses garde-fous côté
sérialiseur :

\begin{itemize}
  \item \texttt{delivery} — entrée de stock (positive uniquement).
  \item \texttt{consumption} — sortie de stock (négative uniquement).
  \item \texttt{adjustment} — correction d'inventaire (signe libre).
  \item \texttt{expiry} — péremption ou casse (signe libre).
\end{itemize}

Un mouvement qui ferait passer le stock résultant en négatif est
refusé, sauf si la raison est un \texttt{adjustment} explicite — ce
qui est cohérent avec une réalité de terrain (compter le stock
physique, corriger l'écart).

% À INSÉRER : capture de l'écran Inventaire avec quelques articles
% (Tubes EDTA, réactifs de glycémie...) et un mouvement en cours
% d'enregistrement dans la modal.
\placeholderFigure{Écran Inventaire}%
{inventory-screen}{capture écran — Inventory.jsx avec liste filtrable
par catégorie, bouton « stock bas seulement » et modal d'enregistrement
d'un mouvement}

\section*{Conclusion du chapitre}
La réalisation a confirmé la pertinence des décisions prises au
chapitre précédent. Le \texttt{BaseModel} a évité des centaines de
lignes répétitives. La séparation stricte \emph{lecture / écriture}
des sérialiseurs DRF a rendu les vues lisibles. Côté frontend,
l'investissement dans un \texttt{Modal} réutilisable et une
bibliothèque d'icônes inline a porté ses fruits dès le troisième
écran. Le chapitre suivant aborde la validation et présente les
résultats obtenus.

\newpage

% =====================================================================
%                CHAPITRE 4 — VALIDATION ET RÉSULTATS
% =====================================================================
\chapter{Validation et résultats}

\noindent
Ce chapitre présente les démarches de validation menées sur la
plateforme : tests fonctionnels, scénarios de bout en bout, et
mesures de performances. Il commente aussi quelques captures qui
illustrent la plateforme en contexte d'usage.

\section{Démarche de tests}

\subsection{Tests unitaires côté backend}
Le backend embarque une couche de tests unitaires couvrant les règles
métier les plus sensibles : calcul du split CNAM, transitions de
statut d'un \texttt{TestOrder}, atomicité des mouvements
d'inventaire, archivage de la valeur d'origine lors d'une correction
biologiste. Les tests s'appuient sur la \texttt{TestCase} de Django
avec une base PostgreSQL en mémoire (\texttt{:memory:}) pour des
exécutions rapides.

\subsection{Tests manuels de bout en bout}
Pour les fonctionnalités transverses au backend et au frontend, j'ai
écrit une liste de scénarios manuels, joués à chaque release : par
exemple, le scénario « walk-in avec CNAM 80 \% » consiste à créer
deux nouveaux patients, sélectionner trois tests, vérifier le split
affiché dans la modal, créer la commande, vérifier la facture
résultante.

\subsection{Vérification de la compilation}
Le pipeline de build local enchaîne :

\begin{enumerate}
  \item Vérification syntaxique de tous les modules Python via
        \texttt{python -m py\_compile}.
  \item Application des migrations sur une base de test temporaire.
  \item Build Vite du frontend web et vérification de l'absence
        d'avertissement.
\end{enumerate}

% À INSÉRER : capture du terminal montrant le build Vite réussi et
% la sortie des tests Django.
\placeholderFigure{Build local — sortie du terminal}%
{build-terminal}{capture écran — terminal montrant le build Vite
réussi (40 modules transformed) et la sortie d'une session de tests
Django}

\section{Scénarios fonctionnels de bout en bout}

\subsection{Scénario 1 : analyse standard avec CNAM}
\textit{Profil patient :} Mme A., 47 ans, couverte CNAM à 80 \%.

\begin{enumerate}
  \item La secrétaire ouvre la modal Walk-in, saisit le numéro de
        Mme A. (le système retrouve son profil et pré-remplit le
        numéro CNAM), sélectionne « NFS » et « Glycémie à jeun ».
  \item Le walk-in crée le rendez-vous, l'échantillon et les deux
        ordres. Le total est de 2~400 MRU ; CNAM 1~920 MRU ; reste à
        charge 480 MRU.
  \item Le technicien réceptionne l'échantillon ; saisit les
        résultats.
  \item Le biologiste valide sans modification. Les ordres passent en
        \texttt{validated}.
  \item La facture est éditée et imprimée pour Mme A.
\end{enumerate}

% À INSÉRER : capture finale de la facture du scénario 1.
\placeholderFigure{Facture finale du scénario 1}%
{scenario1-invoice}{capture écran — facture du scénario 1 (NFS +
glycémie) avec le split CNAM 80 \%}

\subsection{Scénario 2 : correction biologiste tracée}
\textit{Profil patient :} M. B., 35 ans, sans couverture CNAM.

Le technicien saisit la TSH à 2~mUI/L \emph{flag normal}. Le
biologiste, après relecture, requalifie en 2.4 mUI/L \emph{flag haut}
et ajoute un commentaire. Le résultat publié au patient contient la
valeur corrigée ; la base conserve la valeur originale dans
\texttt{original\_value} pour traçabilité.

\section{Mesures de performance}

\subsection{Endpoint statistiques}
La requête \texttt{GET /lab/stats/?days=30} sur un jeu de données de
test contenant 3~120 \texttt{TestOrder} s'exécute en moyenne en
\textbf{180 ms} (mesure côté serveur, médiane sur 50 requêtes,
PostgreSQL local). Le découpage par jour et par catégorie reste
indolore tant que la fenêtre ne dépasse pas 90 jours — borne
choisie pour éviter les agrégations trop longues côté base.

\subsection{Build du frontend}
Le frontend web (40 modules transformés) se construit en moins de
500 ms en moyenne avec Vite. Le bundle final pèse 325 ko gzippé à
91 ko, ce qui reste largement compatible avec une connexion 3G
mauritanienne courante.

\section{Captures en contexte d'usage}

% À INSÉRER : ici, idéalement, 4 à 6 captures montrant l'application
% en utilisation réelle dans un labo. Si vous avez la possibilité de
% capturer des images d'un poste informatique en environnement
% laboratoire (avec autorisations), c'est l'occasion. Sinon, des
% captures plein écran de l'application suffisent.
\placeholderFigure[0.92]{Tableau de bord — vue chef de laboratoire}%
{capture-dashboard}{capture écran — Dashboard.jsx avec KPI, file
d'attente et staff en poste, vue d'un chef de laboratoire}

\placeholderFigure[0.92]{Écran Analyses — onglet « À valider »}%
{capture-analyses-validate}{capture écran — Analyses.jsx onglet
« À valider » côté biologiste, avec quelques ordres en attente}

\placeholderFigure[0.92]{Application mobile — écran d'accueil}%
{capture-mobile-home}{capture écran — application mobile patient,
écran d'accueil avec recherche de laboratoires sur la carte}

\section*{Conclusion du chapitre}
Les tests réalisés confirment que les principales règles métier sont
correctement implémentées et que les performances restent dans des
plages acceptables pour le contexte cible. Les captures, à compléter
au moment du déploiement réel, illustrent que l'expérience
utilisateur respecte le découpage métier et reste lisible sur des
configurations matérielles modestes.

\newpage

% =====================================================================
%                          CONCLUSION GÉNÉRALE
% =====================================================================
\chapter*{Conclusion générale}
\addcontentsline{toc}{chapter}{Conclusion générale}
\markboth{CONCLUSION GÉNÉRALE}{}

\section*{Rappel de la problématique}
Le projet est parti d'une question concrète : peut-on concevoir une
plateforme de gestion de laboratoire qui parte des contraintes
réelles du contexte mauritanien, qui respecte le découpage métier des
équipes sur place, et qui reste accessible à une petite structure ?

\section*{Synthèse des travaux}
La réponse construite, \textbf{labConnect}, se présente comme une
plateforme à trois étages — backend Django, application web pour le
personnel, application mobile pour le patient et le livreur — qui
couvre l'ensemble du cycle de vie d'une analyse, de l'enregistrement
au comptoir jusqu'à la publication du résultat validé et la
facturation tenant compte de la couverture CNAM. Trois modules
transverses (questionnaire pré-test, statistiques agrégées, gestion
d'inventaire) viennent compléter l'ensemble pour offrir un outil
opérationnel à un petit ou moyen laboratoire.

\section*{Principaux résultats}
Sur le plan fonctionnel, la plateforme couvre les six familles de
besoins identifiées lors des entretiens préalables. Sur le plan
technique, l'architecture retenue privilégie l'auditabilité (chaque
opération critique laisse une trace signée et horodatée) et la
performance (les endpoints d'agrégation s'exécutent en moins de
200 ms sur des jeux de données réalistes). Sur le plan UX, l'interface
adapte automatiquement ses écrans aux permissions de l'utilisateur
connecté, évitant les modes ou bascules manuelles.

\section*{Limites et perspectives}
Plusieurs pistes restent à explorer pour faire évoluer la plateforme
vers un produit déployable en environnement médical réel :

\begin{itemize}
  \item \textbf{Connexion aux automates analytiques.} À ce jour, les
        résultats sont saisis à la main par le technicien.
        L'intégration des protocoles ASTM ou HL7 permettrait une
        importation automatique depuis les automates d'hématologie ou
        de biochimie.
  \item \textbf{Génération de codes-barres et de PDF de bilan.} Le
        code-barres du \texttt{Sample} existe déjà comme chaîne de
        caractères ; il manque sa génération graphique (SVG ou PNG)
        pour impression d'étiquettes, ainsi qu'un export PDF
        professionnel du bilan validé.
  \item \textbf{Prescription médecin → laboratoire.} La modélisation
        d'un acteur \emph{médecin prescripteur} ouvrirait la voie à un
        circuit dématérialisé de bout en bout, du cabinet médical au
        laboratoire.
  \item \textbf{Conformité PHI plus stricte.} Le passage en
        environnement de production réel nécessite l'activation
        effective de la 2FA, le chiffrement au niveau colonne des
        données PHI, et la mise en place d'une politique de
        sauvegarde et de rétention conforme à la réglementation
        locale.
\end{itemize}

Au-delà de ces évolutions techniques, le test grandeur nature dans
un laboratoire mauritanien partenaire reste l'étape la plus
importante pour valider l'adéquation entre la conception et le
terrain. Ce projet a constitué une première étape — une base
fonctionnelle, documentée et auditable, prête à être confrontée à
l'usage réel.

\newpage

% =====================================================================
%                          ANNEXES
% =====================================================================
\appendix
\chapter{Détails techniques complémentaires}

\section{Modèle de données complet}
% À INSÉRER : tableau récapitulatif des tables avec leurs champs
% principaux. Possibilité d'utiliser un screenshot de la sortie de
% `manage.py dbshell` ou d'un outil comme DBeaver.

\section{Cartographie des endpoints REST}
La cartographie complète des endpoints REST se trouve dans le fichier
\texttt{API.md} à la racine du dépôt. Le tableau~\ref{tab:endpoints}
en présente une synthèse par domaine.

\begin{table}[ht]
  \centering
  \caption{Synthèse des endpoints REST par domaine}
  \label{tab:endpoints}
  \small
  \begin{tabularx}{\linewidth}{l X}
    \toprule
    \textbf{Domaine} & \textbf{Endpoints principaux} \\
    \midrule
    Authentification & \texttt{/auth/otp/}, \texttt{/auth/login/staff/}, \texttt{/auth/login/google/}, \texttt{/auth/login/firebase/}, \texttt{/auth/refresh/}, \texttt{/auth/me/} \\
    Laboratoires & \texttt{/lab/config}, \texttt{/lab/mine/}, \texttt{/lab/employees/} \\
    Catalogue & \texttt{/lab/catalog/}, \texttt{/lab/catalog/{uuid}/} \\
    Workflow analyses & \texttt{/lab/samples/}, \texttt{/lab/orders/}, \texttt{/lab/orders/{uuid}/result/}, \texttt{/lab/walk-in/} \\
    Facturation & \texttt{/lab/invoices/{appointment\_uuid}/} \\
    Statistiques & \texttt{/lab/stats/?days=N} \\
    Inventaire & \texttt{/lab/inventory/}, \texttt{/lab/inventory/{uuid}/movement/} \\
    Patient mobile & \texttt{/appointments/}, \texttt{/results/mine/} \\
    \bottomrule
  \end{tabularx}
\end{table}

\section{Déclaration d'usage d'outils d'IA générative}
Conformément au § 4 du Guide ID.145, je déclare avoir utilisé un
assistant d'IA générative (\textbf{Claude}, Anthropic) au cours du
projet. Ses contributions, encadrées et systématiquement revues, ont
porté sur :

\begin{itemize}
  \item L'aide à la rédaction et à la reformulation de certains
        passages du présent rapport.
  \item L'écriture de squelettes de code pour gagner du temps sur
        des patterns répétitifs (sérialiseurs, écrans React
        similaires), systématiquement adaptés, simplifiés et
        validés à la main.
  \item La revue critique de certains choix de conception
        (alternatives, contre-exemples), qui m'a permis de tester
        mes propres décisions avant de les figer.
\end{itemize}

Les éléments de fond — modèle de données, choix architecturaux,
décisions UX, démarche d'analyse des besoins — relèvent de mon
travail personnel. Le code livré a été lu, compris, modifié et,
quand nécessaire, réécrit par moi-même.

\newpage

% =====================================================================
%                          BIBLIOGRAPHIE
% =====================================================================
\chapter*{Bibliographie}
\addcontentsline{toc}{chapter}{Bibliographie}
\markboth{BIBLIOGRAPHIE}{}

% Format conforme au guide : [Label]. Auteurs, "titre", édition ou
% conférence, lieu, date.
\begin{thebibliography}{99}

\bibitem{cnam2024}
  Caisse Nationale d'Assurance Maladie, « Guide de l'assuré », CNAM,
  Nouakchott, Mauritanie, 2024.

\bibitem{tate2014lis}
  J.\ Tate, « Laboratory Information Systems: Concepts and Practice »,
  \emph{Clinical Chemistry and Laboratory Medicine}, vol.~52, n\textsuperscript{o}~4,
  pp.~503--515, 2014.

\bibitem{django2024}
  Django Software Foundation, \emph{Django 5.0 Documentation},
  \url{https://docs.djangoproject.com/en/5.0/}, consulté en 2026.

\bibitem{drf2024}
  T.\ Christie, \emph{Django REST Framework Documentation},
  \url{https://www.django-rest-framework.org/}, consulté en 2026.

\bibitem{react2024}
  Meta Platforms, \emph{React Documentation}, \url{https://react.dev/},
  consulté en 2026.

\bibitem{expo2024}
  Expo, \emph{Expo SDK 54 Documentation},
  \url{https://docs.expo.dev/versions/v54.0.0/}, consulté en 2026.

\bibitem{rfc7519}
  M.\ Jones, J.\ Bradley, et N.\ Sakimura, « JSON Web Token (JWT) »,
  IETF, RFC~7519, mai 2015.

\bibitem{firebase2024}
  Google LLC, \emph{Firebase Authentication Documentation},
  \url{https://firebase.google.com/docs/auth}, consulté en 2026.

\bibitem{postgis2024}
  PostGIS Development Group, \emph{PostGIS 3.4 Manual},
  \url{https://postgis.net/documentation/}, consulté en 2026.

\bibitem{nielsen1994}
  J.\ Nielsen, \emph{Usability Engineering}, Morgan Kaufmann, San
  Francisco, 1994.

\end{thebibliography}

\end{document}
